\hwhat{An agent's style vector (17 formatting and punctuation rates) is scored at natural experiments against the agent's own placebo transitions. Round 1 refuted strict conservation: goal switches, forced erasures and assigned registers move style. Round 2 asked whether those moves are speech-act genre, whether they follow an Ornstein--Uhlenbeck drift that an erasure resets, and whether function words are a better charge. \textbf{The moves are not genre.} With 23 genre covariates and context position removed, a forced erasure still moves style ($T=0.546\,[0.522,0.564]$) and not content (0.51, two embedding models). Judges and the Prankster also keep their shifts. \textbf{The OU drift fails}: style variance does not grow with position, and the context-held covariance does not decay. A per-segment offset that the erasure redraws fits (post hoc). Function words move more at goal switches (0.73) but add identity. The holdout is not run.}

\hmodel{Vector spins (model 11). Message $k$ of agent $i$ in a context segment carries the substrate charge $q_i$, day jitter $\eta$ and a context-held term $e_k$. Round 2 pre-registered an OU excursion from zero at each reset. The data favour a constant offset $u_s$, drawn when segment $s$ opens. $\Delta C(l)$ compares style cross-products of same-day message pairs at lag $l$ within a segment and across a forced erasure, matched on time gap. $\rho$ is the pull toward the agent's own last $\le3$ messages.}

\hmath{\vspace{-6pt}\begin{gather*}
x_k=q_i+\eta_{i,d}+e_k+\varepsilon_k,\qquad
\underbrace{e_k=\phi\,e_{k-1}+\sigma\xi_k,\ e_0=0}_{\text{pre-registered (fails)}}\quad\text{vs}\quad \underbrace{e_k=u_s}_{\text{post hoc}}\\
\Delta C(l)=\bigl\langle\!\langle \tilde x_j,\tilde x_{j+l}\rangle\!\bigr\rangle_{\rm within}-\bigl\langle\!\langle \tilde x_j,\tilde x_{j+l}\rangle\!\bigr\rangle_{\rm across},\qquad
\rho=\frac{\sum\langle \tilde x_k,\bar r_k\rangle}{\sum\lVert\bar r_k\rVert^2}\\
T=\Bigl\langle \Pr_{\rm placebo}\!\bigl[D_p<D_b\bigr]\Bigr\rangle,\qquad D=\lVert x-x'\rVert^2/\mathrm{med}_{i,u}(D_{\rm within})
\end{gather*}\vspace{-8pt}}

\hobs{../figures/r2_obs.pdf}{(a) Boundary percentile $T$ by channel (½ = no move; grey band ±0.05; 95\% CIs). Goal switches move every channel, content most. Forced erasures and voluntary consolidations move style, with or without genre and position control. Function words move less, content not at all. (b) The context-held state. $\Delta C(l)$ stays positive and flat up to lag 8 (orange; shaded 95\% CI). A synthetic per-segment offset (squares) has this shape. A synthetic OU drift from zero at the reset (triangles) decays. Both references are scaled to the real lag-1 value.}

\htelos{To attribute messages to a model, combine 50 function-word rates with the 17 formatting rates. Across a goal switch they identify the agent at 0.85, against 0.51 for content embeddings (chance 0.15). After a forced context erasure, expect about 5\% of an agent's style variance to be redrawn, whatever its speech act. Compare an agent's messages within one context segment, or against its mean over many segments. Do not read a style jump at an erasure as a model swap. A judging role or a prankster persona shifts style beyond genre too.}

\hbottom{Style is a stable identity charge plus a context-held offset that each forced erasure redraws (erasure $T$ 0.546 after genre control; content flat). Style plus function words identify agents at 0.85.}

\hresults{\scriptsize\setlength{\tabcolsep}{2pt}\begin{tabular}{@{}p{0.33\columnwidth}p{0.50\columnwidth}p{0.15\columnwidth}@{}}
\textbf{Prediction (dated)} & \textbf{Observed} & \textbf{Verdict}\\\hline
R1-P1 erasure survives genre & $T$ .546 [.522,.564]; genre-matched .545; content .51 & pass\\
R1-P2 \#12 judge register & agent 9 .80 (content .60); one-time judges 4/4 at 1.00 & pass\\
R1-P3 Prankster & percentile 1.00 (5.3$\times$ placebo median) & pass\\
R1-P4 goal shift not genre & $T$ .687 (round 1 .694); identity .78 & pass\\
R2-P1 OU variance growth & $\bar\Delta$ .25 [$-$.16,.59]; $\tau$ not identified & fail (kill)\\
R2-P2 reset, OU decay & $\Delta C(1)$ .70 [.26,1.11]; no decay to lag 8 & mixed\\
R2-P3 self-pull held by context & $\rho$ .36 vs .22 erased; diff .145 [.086,.184] & partial\\
R3-P1 function words conserved & $T_{fw}$ .727 vs style .695 & fail\\
R3-P2 function words at erasure & .528 [.501,.553] & pass (marg.)\\
R3-P3 function-word identity & .79 (style .77); with style .85 & pass\\
\end{tabular}}

\hobsb{../figures/r2_obsb.pdf}{(a) Identity accuracy across 24 goal switches. Each dot is one switch, trained on the last days before and tested on the first days after. Bars are means. (b) Pull toward the agent's own last messages while they are in the context (orange) and after a forced erasure removes them (grey), pooled and per lab (number of agents). DeepSeek has two agents, so its interval is not meaningful.}

\hcaveats{The genre labels are DQ2 reply and stance, mentions and DQ3 window states. They are coarse, and a finer speech-act label might absorb more. Round 1's erasure estimate was inflated by agent-free strata. Scaling by each agent's spread lowers it from 0.564 to 0.547 (A5). The reset-and-hold shape and the lab split are post hoc. Lab groups have 2--11 agents. The erasure effect is strong in \#36--\#42 and weak in \#44 and \#51. The function-word shift is mostly pronouns, so it is partly topic. The holdout is not run.}

\hnext{Test the reset-and-hold offset on reserved data: $\Delta C$ flat in lag, present at the first message after a reset. Find what sets the offset: the memory file reloaded at the reset, or the first messages read. Build the attribution monitor (H46-R4) on function words plus style, with a context-segment correction.}
